\section{The Newton law of an admissible observation germ}
\label{sec:germ-newton-v94}
The order of a determinant perturbation must be computed after its
composition with the admissible model. This section gives such a
construction without assuming that the observation map is an immersion.
The construction uses the monic spectral polynomial. It therefore
continues to make sense when the observed matrix loses rank and its
determinant is identically zero.

Let $X$ be a compact semialgebraic set, let $F:X\to\mathcal P$ be a
continuous semialgebraic map into a finite product of probability
simplices, and let $A_x(z)$ be a continuous semialgebraic family of
monic polynomials of fixed degree $n$. Complex coefficients are
identified with their real and imaginary parts. Put
\[
 h_w(Q,Q')^2=\sum_jw_j\sum_e(\sqrt{Q_{je}}-\sqrt{Q'_{je}})^2,
 \qquad w_j>0,\quad \sum_jw_j=1.
\]
Fix $P\in F(X)$ and write $B_P(t)=\{x\in X:h_w(F(x),P)\le t\}$.
The target is $Z(A_x)$, including multiplicity. Its metric is
\[
 d_\infty(Z,Z')=\min_{\pi\in\mathfrak S_n}\max_i|z_i-z'_{\pi(i)}|,
 \qquad |\cdot|\text{ is the Euclidean modulus on }\C.
\]
Real polynomials impose conjugacy constraints; those constraints are
part of $X$ and are never discarded. Define
\begin{equation}\label{eq:germ-modulus-v94}
 \omega_{P,X}(t)=\diam\{Z(A_x):x\in B_P(t)\}.
\end{equation}
The whole inverse image of the observation ball is used. A pointed
version is obtained by specifying a compact local admissible set $X$
instead; it is not identified with the whole-model version.

Assume that the spectral polynomial is constant on the exact fibre:
\begin{equation}\label{eq:spectral-fibre-v94}
 A_x=A_0\quad\hbox{for every }x\in F^{-1}(P).
\end{equation}
This is identification of the target, not of the latent parameter.
Write $A_0(z)=\prod_{r\in\mathcal R}(z-r)^{m_r}$.
Compactness and \eqref{eq:spectral-fibre-v94} imply that $A_x\to A_0$
uniformly for $x\in B_P(t)$ as $t\downarrow0$. In disjoint discs
about the distinct $r$, the corresponding monic cluster factors are
therefore uniquely defined for small $t$:
\begin{equation}\label{eq:cluster-coeff-v94}
 A_{r,x}(r+\zeta)=\zeta^{m_r}
                   +\sum_{a=0}^{m_r-1}c_{r,a}(x)\zeta^a,
 \qquad A_x=\prod_r A_{r,x}.
\end{equation}
These coefficients are analytic functions of the coefficients of $A_x$
in a neighbourhood of $A_0$, and are semialgebraic there. Indeed,
the differential of multiplication of the coprime monic cluster
factors is invertible: reducing a zero differential modulo any one
factor forces its own variation to vanish. The analytic inverse
function theorem and the locally unique polynomial factor equations
prove the two assertions.

For each of the finitely many pairs $(r,a)$ put
\begin{equation}\label{eq:coefficient-envelope-v94}
 b_{r,a}(t)=\max_{x\in B_P(t)}|c_{r,a}(x)|,
 \qquad \beta_{r,a}=\operatorname{ord}_{t=0}b_{r,a}(t).
\end{equation}
The order is $+\infty$ when the envelope vanishes on a neighbourhood
of zero. Otherwise it is the positive exponent of its leading
Puiseux monomial. Define
\begin{equation}\label{eq:germ-edge-v94}
 \alpha_P=\min_{r,a<m_r}\frac{\beta_{r,a}}{m_r-a}.
\end{equation}
The order in \eqref{eq:coefficient-envelope-v94} is an observation
order. It includes all orders of the observation map, the spectral
map, and the admissibility constraints. It is not the order of a
coefficient on an affine tangent space.

Suppose $\alpha_P<\infty$. Consider the joint, rather than
coordinatewise, rescaled coefficient sets
\begin{equation}\label{eq:scaled-coefficients-v94}
 \mathcal C_t=
 \left\{\left(t^{-\alpha_P(m_r-a)}c_{r,a}(x)\right)_{r,a}:
                                           x\in B_P(t)\right\}.
\end{equation}
Let $\mathcal C_0$ be the fibre at $t=0$ of the closure of their
graph for $0<t<t_0$. For $v\in\mathcal C_0$ let
\[
 \mathcal Z_r(v)=Z\left(\zeta^{m_r}
                              +\sum_{a<m_r}v_{r,a}\zeta^a\right),
 \quad
 C_P=\max_{v,v'\in\mathcal C_0}\max_r
                         d_\infty(\mathcal Z_r(v),\mathcal Z_r(v')).
\]
The common parameter in \eqref{eq:scaled-coefficients-v94} retains
all cross-cluster, conjugacy, and stochastic constraints. Replacing
$\mathcal C_0$ by a product of its coordinate projections would in
general give an incorrect leading constant.

\begin{theorem}[Intrinsic Newton law]\label{thm:germ-newton-v94}
Under the preceding hypotheses, either $\alpha_P=\infty$ and the
spectral target is constant on $B_P(t_0)$ for some $t_0>0$, or
$\alpha_P\in\mathbb Q_{>0}$, $0<C_P<\infty$, and
\begin{equation}\label{eq:germ-exact-v94}
 \omega_{P,X}(t)=C_Pt^{\alpha_P}+o(t^{\alpha_P}).
\end{equation}
Moreover $\mathcal C_t\to\mathcal C_0$ in Hausdorff distance.
The exponent and the leading root-multiset set depend only on the
germ of the observation--spectral image $(F,A)(X)$ at $(P,A_0)$,
with its specified observation metric. They are unchanged by
analytic or semialgebraic reparametrizations preserving that image
germ, including its admissibility constraints.

For real algebraic input, the exponents, the limiting coefficient
set, and its variational diameter admit a finite description by
real quantifier elimination and algebraic branch analysis. No
polynomial-time claim or model-independent jet cutoff is made.
In particular the theorem applies to the entire closed stochastic
model \eqref{eq:model}, with $A_\theta=\prod_b f_b$, without a
channel-rank assumption, without $\tau(P)>0$, and without positive
observed-cell assumptions, whenever its exact spectral fibre is a
singleton.
\end{theorem}
\begin{proof}
All maxima in \eqref{eq:coefficient-envelope-v94} exist. The graph
of each is semialgebraic: introduce nonnegative square-root variables
for the finitely many cell probabilities, use the polynomial equations
for the cluster factors, and quantify the condition of being a maximum.
Each envelope tends to zero. One-variable semialgebraic preparation
therefore gives either local zero or a positive leading Puiseux
monomial. These are the standard real-algebraic ingredients
\cite{BCR98,BPR06}; the remaining argument identifies what their
orders and leading sets mean for the observed spectrum.

For a monic polynomial $p(\zeta)=\zeta^m+\sum_{a<m}u_a\zeta^a$,
let $R(p)$ be its largest root modulus and put
$H(p)=\max_{a<m}|u_a|^{1/(m-a)}$. Elementary symmetric functions
and a geometric-series root bound give constants depending only on
$m$ such that
\begin{equation}\label{eq:root-coeff-comparison-v94}
 c_m H(p)\le R(p)\le2H(p).
\end{equation}
For the first inequality use
$|u_a|\le\binom ma R(p)^{m-a}$. For the second, if
$|\zeta|>2H(p)$, then
$\sum_{a<m}|u_a\zeta^a|<|\zeta|^m$.
Thus the maximal displacement from the base spectrum has order
\eqref{eq:germ-edge-v94}. Its diameter is between that maximal
displacement and twice it, since the base spectrum belongs to every
ball. This proves the exponent, but not yet the exact constant.

By the definition of $\alpha_P$, every coordinate of
\eqref{eq:scaled-coefficients-v94} is uniformly bounded. The sets
$\mathcal C_t$ are nonempty compact semialgebraic sets in a fixed
compact box. Their set of cluster limits $\mathcal C_0$ is nonempty,
compact, and semialgebraic. We verify Hausdorff convergence, rather
than replacing a subsequential limit by a limit without justification.
The outer inclusion follows from compactness: any sequence violating
it has a convergent subsequence belonging to $\mathcal C_0$.
For fixed $v\in\mathcal C_0$, the bounded semialgebraic function
$t\mapsto\operatorname{dist}(v,\mathcal C_t)$ has a limit at zero.
Its definition as a cluster point supplies a subsequence along which
this distance tends to zero, so the limit is zero. A finite net in
$\mathcal C_0$, and the fact that distance to a set is 1-Lipschitz,
make this convergence uniform in $v$. This proves the inner inclusion.

After the exact change $z=r+t^{\alpha_P}\zeta$, the monic cluster
polynomial divided by $t^{\alpha_Pm_r}$ has precisely the
coefficients in \eqref{eq:scaled-coefficients-v94}. Root multisets
are continuous functions of monic coefficients in the bottleneck
metric, uniformly on compact coefficient sets. Hence their joint
rescaled sets converge to $(\mathcal Z_r(v))_r$, $v\in\mathcal C_0$.
Small disjoint base discs prevent a minimizing matching from crossing
between different base clusters. Taking diameters proves
\eqref{eq:germ-exact-v94}. The set $\mathcal C_0$ contains zero.
At least one minimizing envelope has a positive leading coefficient,
so it also contains a nonzero vector. The corresponding collection
of monic polynomials cannot all be pure powers. Its distance from
the zero-root collection is positive. Thus $C_P>0$.

If all orders are infinite, all finitely many cluster coefficients
vanish on a common observation neighbourhood. This proves the
rigidity alternative. Conversely, a constant spectral target makes
all these coefficients zero. The entire construction takes place
in $(F,A)(X)$, proving the asserted germ invariance, not merely basis
invariance of a tangent space.

For algebraic input, factor isolation, the graphs of the envelopes,
and the closure defining $\mathcal C_0$ have finite formulas over
real algebraic numbers. Quantifier elimination constructs their
semialgebraic descriptions, and Puiseux branch analysis gives the
orders. Polynomial root equations and the finitely many matchings
then describe $C_P$ by a further compact optimization formula.
These are effective finite procedures, not claims of efficient
computation. In the stochastic model, compactness, semialgebraicity,
and positivity of the clock denominators were established in the
model definition; $\prod_bf_b$ remains monic even when $K$ loses
rank. All hypotheses are therefore verified as stated.
\end{proof}

\subsection{Composed analytic jets}
The preceding theorem handles singular observation maps and inequality
faces directly. When an admissible germ is an immersed analytic chart,
its coefficient envelopes can instead be calculated from ordinary
homogeneous jets. The following statement makes the required
composition explicit.

Let $x\mapsto M(x,z)$ be real analytic for $x$ in a neighbourhood
of zero in $\R^N$, polynomial of fixed degree in $z$, with invertible
leading matrix. Suppose its probability-normalized observation map
$F(x)$ has positive cells and injective derivative $J=DF(0)$.
Choose $\rho>0$ so small that $F$ is an embedding on a neighbourhood
of the closed ball $\|x\|\le\rho$. All competitors in this subsection
are explicitly restricted to that ball. Put
\[
 I(v)=\sum_jw_j\sum_e\frac{(Jv)_{je}^2}{P_{je}},
 \qquad \mathcal E=\{v:I(v)\le4\}.
\]
At a root $r$ of $D_0(z)=\det M(0,z)$ of multiplicity $m_r$ write
\begin{equation}\label{eq:composed-jets-v94}
 \det M(x,r+\zeta)=D_0(r+\zeta)
            +\sum_{j\ge1}\sum_a H_{r,j,a}(x)\zeta^a,
 \quad H_{r,j,a}\text{ homogeneous of degree }j.
\end{equation}
The determinant in \eqref{eq:composed-jets-v94} is taken after
substitution of the complete map $M(x)$.
Set
\[
 \alpha=\min_{r,j,a<m_r:H_{r,j,a}\not\equiv0}\frac{j}{m_r-a},
 \quad
 \Pi_{r,v}(\zeta)=b_r\zeta^{m_r}
       +\sum_{j=\alpha(m_r-a)}H_{r,j,a}(v)\zeta^a,
 \quad b_r=\frac{D_0^{(m_r)}(r)}{m_r!}.
\]
Empty sums are zero; an empty minimum is infinite.

\begin{theorem}[Composed-jet formula]\label{thm:composed-jets-v94}
If $\alpha<\infty$, the local observation-ball diameter equals
\[
 t^\alpha\max_{v,v'\in\mathcal E}\max_r
       d_\infty(Z(\Pi_{r,v}),Z(\Pi_{r,v'}))+o(t^\alpha).
\]
The displayed constant is finite and positive. Once any nonzero
coefficient below a base multiplicity has supplied a provisional
upper bound $\alpha_0$ on $\alpha$, only parameter jets of orders
$j\le\lfloor\alpha_0\max_r m_r\rfloor$ can lower that bound or
contribute to a minimizing edge. The all-zero case is local spectral
rigidity and, for a general analytic germ, requires an identity
argument rather than any fixed finite test of Taylor coefficients.
The formula is covariant under analytic changes of parameter with
invertible derivative, provided the complete admissible and observed
maps are composed with that change.
\end{theorem}
\begin{proof}
The square-root expansion gives
$h_w(F(tv),P)=\tfrac12t\sqrt{I(v)}+O(t^2)$ uniformly on bounded
sets. Injectivity of $J$ bounds $\|x\|$ by a constant times the
observation distance in the chosen neighbourhood. The rescaled
parameter balls converge to $\mathcal E$, including its boundary
by radial approximation.

For $x=tv$ and $z=r+t^\alpha\zeta$, division by
$t^{\alpha m_r}$ leaves precisely $\Pi_{r,v}$ in the limit.
The convergence is uniform on bounded $v,\zeta$ sets. To justify
this assertion for an infinite analytic expansion, truncate at an
integer $J_0>\alpha_0\max_r m_r$. Analyticity bounds the remaining
coefficient functions by $O(t^{J_0+1})$ uniformly on a smaller
parameter ball; after the division their contribution tends to
zero. All retained nonedge terms have strictly larger weight.
On $|\zeta|=R$, first choosing $R$ sufficiently large and then
$t$ small makes the perturbation smaller than the base leading
monomial. Rouch\'e's theorem, applied also on small fixed circles
around the base clusters, localizes exactly $m_r$ roots and excludes
loss at infinity. Continuity of root multisets now gives uniform
convergence, and taking diameters proves the formula. A nonzero
edge polynomial at an interior point of $\mathcal E$ proves
positivity by comparison with $v=0$.

Terms of degree $j>\alpha_0\max_r m_r$ have
$j/(m_r-a)>\alpha_0$ and cannot contribute. If every lower
coefficient vanishes identically, each determinant retains all
base roots with their full multiplicities and the degree is
unchanged; it is a nonzero scalar multiple of the base determinant.
Finally, an analytic parameter diffeomorphism is a bijection of
local admissible alternatives. It transports the observation balls
exactly, and its derivative transports their limiting ellipsoids.
Equivalently, the lowest nonzero homogeneous part of each composed
coefficient transforms by that derivative. This proves the stated
covariance while retaining, rather than discarding, higher model jets.
\end{proof}

For $M(x)=M_0+E_1(x)+E_2(x)+\cdots$, its quadratic determinant
jet is
\[
 D\det_{M_0}[E_2(x)]
       +\tfrac12D^2\det_{M_0}[E_1(x),E_1(x)].
\]
Both summands are necessary. The affine result
Theorem~\ref{thm:constrained-newton} is the specialization in which
$E_j=0$ for $j\ge2$. For a singular observation derivative one must
use Theorem~\ref{thm:germ-newton-v94}, not a degenerate Fisher
ellipsoid in the last theorem.

\subsection{A curved cancellation with the same affine tangent}
\begin{proposition}[Arbitrarily delayed composed exposure]
\label{prop:curved-cancellation-v94}
For every $m\ge3$ and integer $q>2$, a positive rational polynomial
experiment has an analytic one-parameter admissible germ whose
affine tangent predicts exponent $2/m$, whereas its actual exponent
is $q/m$. A corresponding germ with the same affine tangent can
also be spectrally constant.
\end{proposition}
\begin{proof}
Take $f=z^m$ and $g=\prod_{i=1}^m(z-a_i)$ with distinct $a_i>0$,
and at least $2m+1$ clocks greater than all $a_i$. Let
$O=2^{-1/2}\left(\begin{smallmatrix}1&1\\1&-1\end{smallmatrix}\right)$
and define $M_x=\tfrac12OB_xO^{\mathsf T}$, where
\[
 B_x=\begin{pmatrix}f+x^2+x^q&xg\\x&g\end{pmatrix}.
\]
The entry sum of $M_x$ is $f+x^2+x^q$, monic of degree $m$.
At $x=0$ all observed cells are positive since $0<g(T)<f(T)$;
positivity persists locally. The coprimality argument in
Proposition~\ref{prop:quadratic-exposure} gives $\tau(P)>0$.
The first derivative of the observation is nonzero by
Lemma~\ref{lem:rational-chart-v93} applied to the nonzero tangent.
But the exact identities are
\[
 \det B_x=g(f+x^q),\qquad
 \det(B_0+xB'_0)=g(f-x^2).
\]
At zero these give exponents $q/m$ and $2/m$, respectively.
The roots of $g$ stay fixed; the roots of $f+x^q$ have modulus
$|x|^{q/m}$. Taking the exact local Hellinger diameter therefore
has the claimed exponent even when it exceeds one. Dropping $x^q$
makes the determinant $fg$ identically. These are germs in the
positive rational extension, not asserted stochastic factorizations.
They show why no jet cutoff depending only on matrix size and
polynomial degree can cover arbitrary analytic parametrizations.
\end{proof}

\subsection{Statistical consequence}
\begin{corollary}[The intrinsic local risk exponent]
\label{cor:germ-risk-v94}
In the nonrigid case of Theorem~\ref{thm:germ-newton-v94}, take
independent categorical samples with $n_j/N\to w_j>0$, and fix
$0<c\le1/(8\sqrt2)$. On the anchored observation neighbourhood
$B_P(cN^{-1/2})$, the minimax squared bottleneck risk $R_N$ satisfies
\[
 \frac{3}{32}C_P^2c^{2\alpha_P}
 \le\liminf_{N\to\infty}N^{\alpha_P}R_N
 \le\limsup_{N\to\infty}N^{\alpha_P}R_N
 \le C_P^2c^{2\alpha_P}.
\]
The same assertion holds for the compact local positive rational
experiment of Theorem~\ref{thm:composed-jets-v94}. The diameter
constant is exact; this sandwich is not an assertion of an exact
minimax decision constant.
\end{corollary}
\begin{proof}
The two-point proof of Theorem~\ref{thm:modulus-minimax} uses only
the triangle inequality, the compact target diameter, and categorical
Hellinger product bounds, so applies to both specified experiments.
Substitute \eqref{eq:germ-exact-v94}, or the composed-jet diameter,
into its bounds $3\omega^2/32\le R_N\le\omega^2$.
\end{proof}
